\section{Theoretical Setup}

\subsection{Problem Setting}

State the mathematical setting in paper-ready notation, using only objects from
the formalized setting. Write ordinary LaTeX prose and avoid raw markdown
heading syntax.

\subsection{Assumptions}

State every theorem-critical assumption as a separate numbered assumption
environment. Use the displayed numbering from the theorem environment, so the
paper reads as Assumption 1, Assumption 2, and so on. Do not hide assumptions
in prose or introduce assumptions later in the proof. Preserve the stable
assumption ids from setting.md as labels, for example assump:smoothness becomes
\label{assump:smoothness}.

\begin{assumption}[Short descriptive name]\label{assump:short-name}
State one coherent theorem-critical assumption from the formalized setting.
\end{assumption}

\begin{assumption}[Short descriptive name]\label{assump:second-name}
State the next theorem-critical assumption from the formalized setting.
\end{assumption}

\subsection{Formalized Goal}

State the target quantity and proof objective from the formalized goal in
paper-ready prose.
